\documentclass[11pt]{article}
\usepackage[a4paper,margin=23mm]{geometry}
\usepackage[T1]{fontenc}
\usepackage{lmodern,amsmath,amssymb,amsthm,mathtools,microtype,booktabs}
\usepackage[colorlinks=true,linkcolor=blue,citecolor=blue,urlcolor=blue]{hyperref}
\newcommand{\E}{\mathbb E}
\newcommand{\tr}{\operatorname{tr}}
\newcommand{\Wick}{\mathcal W}
\newcommand{\norm}[1]{\left\lVert#1\right\rVert}
\newtheorem{theorem}{Theorem}
\newtheorem{lemma}{Lemma}
\title{Ordinary charged fields at fixed half filling\\
\large Fourier conditioning, internal contractions and the same-copy strong bridge}
\author{Alan Day \and Chengheng\thanks{AI research assistant (OpenAI).}}
\date{8 November 2026\quad Local research draft}
\begin{document}\maketitle
\begin{abstract}
We extend the central-sector marked-source calculus to ordinary charged
blocks with internal Gaussian contractions. Exact charge projection,
with an integrable Fourier bound after singleton cancellation, subtracts
$2\lambda(\sum_a i_a z_a)^2$ from the full-trace source exponent.
Here $i_a$ counts internal chords in charged block a. This removes the
shared charge fluctuation and establishes the ordinary-to-Wick strong
bridge in every fixed finite central Schatten norm. Exact energy
pinching and fixed polynomial preparation bounds then transfer each
fixed ordinary coefficient to the single-selected state. Uniform tails
of the actual infinite ordinary series remain to be established.
\end{abstract}

\section{Definitions and the missing conditioning step}
Use standard CAR, a Hermitian overlapping complex Gaussian p-body
Hamiltonian H with the same disorder in all fields, even N, m=N/2 and
$\lambda_N=p^2/N\to\lambda>0$. Let $C_Q=\E H^2$, with sector value
$C_m$, and $h_N=H_m/\sqrt{C_m}$. Use normalized annealed sector norms
$\norm{X}_{r,m}=(\E\tr_m|X|^r)^{1/r}$. All degrees and source counts
are fixed before N tends to infinity. Time uses dimensionless h units,
with $\hbar=1$.

The body loss and half-size maps are
\[
 DX=QX-\sum_i c_i^\dagger Xc_i,\qquad
 \mathcal LX=\tfrac12\bigl(NX-\sum_i(c_i^\dagger Xc_i+c_iXc_i^\dagger)\bigr).
\]
They act before sector restriction. Disjoint charged D marks may now
contain complete Gaussian pairs. Neutral L marks still have no internal
pair in their actual marked subinterval. Larger Wick exclusions and
Gaussian-chaos projections retain their actual regions.

For a pairing $\pi$, let $i_a$ be the number of internal chords in
charged mark a, and zero for an L mark. Put
$I(z)=\sum_a i_a z_a$, with commuting $z_a^2=0$.
The old charge-tilt formula on compact complex source domains
\cite{tilt} suggests a conditional rule, but does not itself prove
convergence at the single lattice charge m. Here we perform that
projection explicitly, following the Fourier approach to fixed-charge
complex SYK in \cite{bnr} and retaining the marked-source cancellations.

\newpage
\section{Exact singleton cancellation before Fourier projection}
At one site write $n=c^\dagger c$, $\bar n=1-n$. The symmetric local
body loss gives
\[
 (1-zD)n=(1-z)n,\quad
 (1-zD)\bar n=\bar n+zn,\quad
 (1-zD)c^{(\dagger)}=(1-z/2)c^{(\dagger)}.
\]
The second identity is crucial: an occupation with zero unmarked
weight can acquire a pure source term. Thus an ordinary mark cannot
be treated by deleting the no-internal-pair assertion from the pure
Slater algorithm.

For a Fourier angle $\theta$, use the normalized complex product trace
with $\rho_\theta=e^{i\theta}/(1+e^{i\theta})$. An internal charged
chord has the two singleton ratios
\[
 (1-z_a)(1+e^{i\theta}z_a)=1+(e^{i\theta}-1)z_a.
\]
Each label appears p times. The complete ideal singleton residual is
\begin{equation}
 \exp_{\rm nil}\{p(e^{i\theta}-1)I(z)\}.\label{singleton}
\end{equation}
Open endpoints cancel the $pd_a$ number source exactly; L singleton
deficits cancel in the same way, without an internal-chord term.
The scalar singleton normalization and the unmarked charge trace give
\begin{equation}
 [\cos(\theta/2)]^{N-2kp},\label{cosine}
\end{equation}
for k Gaussian chords. This factor must be kept until the extensive
source terms have cancelled. It tends to a Gaussian after the change
of variable $t=p\theta$.

Let $F_{\pi,N}(\theta,z)$ denote the normalized contribution of this
pairing to the full-Fock trace with
$e^{i\theta(Q-N/2)}$ and the indicated deficit sources. Include
$C_m^{-k}$ and the disorder normalization. Per chord, the combined
label-count and source-free singleton normalization is
\[
 A_N=\frac{4^{-p}\binom Np^2}{\binom mp\binom{m+p}p}
 \longrightarrow e^{-\lambda}.
\]
It is uniformly bounded; the remaining overlap factors are treated
below. For every fixed real t,
the normalized collision expansion gives, coefficientwise,
\begin{equation}
 F_{\pi,N}(t/p,z)\longrightarrow
 e^{-t^2/(8\lambda)}
 \exp_{\rm nil}\{\lambda K_\pi(z)+it I(z)\}.\label{local-limit}
\end{equation}
$K_\pi$ is the existing full-trace source exponent, including
$K_\pi(0)=-4\operatorname{cr}\pi$ and any Gaussian-number shifts.
To obtain \eqref{local-limit}, divide each overlap-site polynomial by
its source-bearing singletons. On compact t intervals their
coefficients tend to their half-occupation values. Fixed label-pair
overlaps have uniformly bounded exponential moments, and triple
memberships vanish after domination. Formula \eqref{singleton} tends
to $e^{itI}$. This extends the previous one/two-source argument to any
fixed source count: the coefficient norm merely acquires a fixed
polynomial power in the total overlap count.

\newpage
\section{An integrable bound for the full Fourier interval}
We give the missing bound that permits conditioning on a single charge.
Let R be the sum of pairwise intersection sizes of the $2k$ independent
p-set labels. For fixed k and every fixed positive a,
\[
 \sup_N\E e^{aR}<\infty.
\]
Indeed each pair overlap Y obeys
$\E e^{aY}\le\exp[\lambda_N(e^a-1)]$ by its factorial moments;
Holder handles their finite sum without assuming independence.

Fix $0<\delta<\pi/2$ and r source variables. On $|\theta|\le\delta$,
the singleton denominators have nonzero constant terms uniformly away
from zero. At most $2R$ label memberships are missing from ideal
singleton sites. After the exact cancellation in \eqref{singleton},
all corrections from those missing sites and overlap-site matrices
have coefficient norm at most
\[
 K(1+R)^r A^R.
\]
This follows by expanding the fixed-r nilpotent products: the local
matrix maps have uniformly bounded coefficients, the number of
non-singleton sites is at most R, and at most r marked factors may be
chosen. No division by the possibly zero unmarked overlap numerator
is used. The residual \eqref{singleton} adds at most
$K(1+p|\theta|)^r$. Consequently each coefficient satisfies
\begin{equation}
 |[z_S]F_{\pi,N}(\theta,z)|
 \le K(1+p|\theta|)^r|\cos(\theta/2)|^{N-2kp}.
 \label{near}
\end{equation}
The overlap average has been absorbed into K. Constants here depend on
the fixed diagram, r and a compact positive range of $\lambda_N$.
For sufficiently large N, $N-2kp\ge N/2$, so under $t=p\theta$,
\begin{equation}
 |[z_S]F_{\pi,N}(t/p,z)|
 \le K(1+|t|)^r e^{-t^2/(16\lambda_+)}.
 \label{domination}
\end{equation}

For $\delta\le|\theta|\le\pi$, avoid complex singleton denominators.
Every label configuration has support U at most $2kp$ sites. Loss
on sites outside that support is zero because $D_iI=\mathcal L_iI=0$.
Sites outside the support therefore contribute exactly a factor
of modulus $|\cos(\theta/2)|^{N-U}$. Direct expansion of the r loss
and number factors gives at most $Kp^r$ on the support. The normalized
label count contributes at most $e^{Kp}$, since k is fixed. Thus
\begin{equation}
 |[z_S]F_{\pi,N}(\theta,z)|
 \le Kp^r e^{Kp}|\cos(\theta/2)|^{N-2kp}.
 \label{far}
\end{equation}
Its integral, even multiplied by p, tends to zero exponentially in N.
Equations \eqref{domination} and \eqref{far} justify dominated
convergence on the rescaled whole Fourier interval. This is stronger
than compact-t convergence of a charge generating function.

\newpage
\section{The central source and the ordinary strong bridge}
Exact half-charge extraction divides the integral of
$F_{\pi,N}(\theta,z)$ by that of $\cos(\theta/2)^N$.
Using \eqref{local-limit}--\eqref{far}, the limiting ratio is
\begin{equation}
 e^{\lambda K_\pi(z)}
 \frac{\int_{\mathbb R}e^{-t^2/(8\lambda)}e^{itI(z)}\,dt}
      {\int_{\mathbb R}e^{-t^2/(8\lambda)}\,dt}
 =\exp_{\rm nil}\{\lambda[K_\pi(z)-2I(z)^2]\}.
 \label{conditioned}
\end{equation}
The Gaussian integral is a finite polynomial identity in the sources.

\begin{theorem}[Fixed central ordinary source]
For every fixed admissible marked layout, the half-filled sector limit
is given by the full-trace source with
\[
 K_{\pi,\mathrm{cen}}=K_\pi-4\sum_{a<b}i_ai_bz_az_b.
\]
All higher fixed central moments exist and are bounded along the
double-scaling sequence. No source correction is added when a mark
has no internal chords. Wick exclusions and chaos projections may
be imposed before summing the finite set of diagrams.
\end{theorem}

For ordinary G sources, the explicit single-source coefficient is
$-4\operatorname{cr}(O_a)-2i_a$, where $O_a$ is the open-chord set.
The connected DD coefficient becomes
$2b(b-1)-4\operatorname{cr}(B)$, with $B=O_1\cap O_2$, $b=|B|$.
The former $4i_1i_2$ term is absent. L has zero internal count, so
the DL and LL formulas retain their previous values.

Normalize the actual ordinary field by the target-sector scalar:
\[
 G_{n,N}=C_m^{-n/2}
 [p\mathcal N-D-\lambda_N\mathcal N(\mathcal N-1)]H^n|_m.
\]
Let $h^n=\sum_d a_{nd}(q)H_d^{(q)}(h)$, $q_N=e^{-4\lambda_N}$,
$m_d(q)=q\partial_q\log[d]_q!$, and define
\begin{align*}
 g_n^{\rm cen}(h)&=\sum_d a_{nd}(q)
     [-4\lambda m_d(q)-\lambda(n-d)]H_d^{(q)}(h),\\
 J_{d,N}&=C_m^{-d/2}(dp-D)\Wick_d(H^d)|_m
 -\lambda_N[d(d-1)-4m_d(q_N)]H_d^{(q_N)}(h_N).
\end{align*}
In the finite-N $g_n$ and $a_{nd}$ below use $q_N,\lambda_N$.
Cutting a fixed internal matching leaves its original $q^{c_\gamma}$
weight. Its remaining source is
$\lambda[-4\operatorname{inv}\sigma-2i]$; its Mahonian average is
$\lambda[-4m_d-2i]$. Subtracting $g_n^{\rm cen}$ thus leaves the
pure J source. The same grouping applies to RR, RJ and JJ Gram terms
on the same original disorder, exactly as in \cite{joint}, now with
the central rule proved independently. Therefore, for fixed n,
\begin{equation}
 \left\|G_{n,N}-g_n^{\rm cen}(h_N)
       -\sum_d a_{nd}(q_N)J_{d,N}\right\|_{2,m}\longrightarrow0.
 \label{bridge}
\end{equation}
All fixed higher moments are bounded by the source theorem. Interpolation
upgrades \eqref{bridge} to every fixed finite Schatten exponent;
Holder permits fixed left/right energy powers. This is a same-copy
strong statement, not equality of two marginal laws.

\newpage
\section{Actual pinching, selected preparations and a scalar calibration}
The counterterm $g_n^{\rm cen}(h_N)$ commutes exactly with $h_N$.
The preceding higher-Wick stage \cite{higher} established
$\norm{\mathcal P_{h_N}J_{d,N}}_{r,m}\to0$ for every fixed d,r,
using the central two-energy measure and its zero diagonal mass.
Apply the contraction of exact pinching to \eqref{bridge}. Then
\begin{equation}
 \left\|(1-\mathcal P_{h_N})G_{n,N}
       -\sum_d a_{nd}(q_N)J_{d,N}\right\|_{r,m}\longrightarrow0
 \label{actual}
\end{equation}
for every fixed n and finite r. Complete degenerate energy blocks are
retained; no long-time averaging limit is interchanged.

For an existing fixed-degree stabilizer-covariant preparation $P_NO$
with bounded norm, the orbit-density estimate \cite{prep} gives
\[
 \E\norm{X_NP_NO}^2\le C_P\norm{X_N}_{4,m}^2.
\]
Thus \eqref{actual} holds on such selected preparation vectors.
Polarization, higher moments and trace Holder transfer any fixed
covariant block between polynomial preparations. Fixed real h-time
phases are handled by the already summable Gaussian energy Taylor
tails. In particular, all fixed ordinary charged truncations now have
an actual, exactly pinched selected-state limit. The limit uses the
charged Wick tower's density $f(h)=1-\sqrt q+\sqrt q\,h^2$.

An exact finite-N check exposes the role of internal contractions.
For m=N/2 the covariance ratio is
\[
 \frac{C_{m-1}}{C_m}
 =\frac{(m-p)(m+p+1)}{m(m+1)},\qquad
 \frac{\E G_{2,N}^{\rm unnormalized}}{C_m}
 =m\left(\frac{C_{m-1}}{C_m}-1\right)
 =-\frac{p(p+1)}{m+1}=:b_{N,p}.
\]
This tends to $-2\lambda$. At n=2 the full difference in
\eqref{bridge} is exactly the deterministic scalar
$(b_{N,p}+2\lambda_N)I$, since $Z_2|_m/C_m=h_N^2-1$.
For $N=4p^2$,
\[
 b_{N,p}+2\lambda_N=\frac{1-2p}{4p^2+2}\longrightarrow0.
\]
Two separately internal G2 blocks have a central diagram weight
$4\lambda^2$. Their full-trace weight is $4\lambda^2+4\lambda$.
The difference is the shared charge variance removed by
\eqref{conditioned}; deleting it without a charge-conditioning
argument would have been unjustified.

This closes the fixed-degree ordinary strong bridge. The full selected
bilocal also needs a uniform finite-N bound on the tail as n grows.
The constants in \eqref{near} depend on the fixed diagram and cannot
be used as such a bound. The known full-Fock Clifford Riesz estimate
\cite{belp} does not automatically apply to fixed-charge trace.
Neither the Fourier argument nor preparation-density bounds authorize
interchanging the infinite ordinary sum with the N limit.

\newpage
\section{Finite evidence and continuation}
Three fresh outputs support separate parts of the construction.
The exact algebra checker verifies 2,186 nilpotent Fourier transforms,
20 energy regressions, 96 RR/RJ/JJ Gram equalities across 32 layouts,
two nonzero full-versus-central G2 gaps, and five finite-N G2 scalar
bridges. It also saves 30 exact conditional-singleton coefficients
at $N=4p^2$, $p=16,32,64,128,256$, including their finite-size errors.

The first checker had a coverage defect: its requested two-mark
local loop generated no cases. That output is retained and supplies
no local-source coverage. A separate checker explicitly enumerates
the intervals and requires positive coverage. It verifies 3,496
two-source, 459 three-source and 1,265 four-source cases against the
unrestricted-site matrix source engine, preserving 1,380 cases with
internal pairs and 4,010 negative nonconstant coefficients.

For disjoint singleton supports, set $L=N-2kp$ and take r internal
source counts $i_1,\ldots,i_r$. The exact normalized half-charge
coefficient is
\[
 \frac{p^r\prod_a i_a}{\binom L{L/2}}
 \sum_{j=0}^r(-1)^{r-j}\binom rj\binom L{L/2-j}.
\]
It tends to zero for odd r and to
$(-4\lambda)^{r/2}(r-1)!!\prod_a i_a$ for even r. This isolates the
conditioning factor; it is not the complete interacting diagram with
overlap corrections. An independent Fourier quadrature checks ten
finite systems against these exact rational values at 60 and 90
decimal digits. A Gaussian incomplete-integral bound controls the
omitted $|t|>16$ interval. Quadrature and roundoff are assessed by
precision comparison, not interval certification.

The next task is the central finite-N ordinary degree tail. Sectorwise
pure-Wick tagged bounds and covariance CP estimates remain available;
the Riesz-dependent product step needs replacement or another uniform
argument. The existing full-trace bilocal four-point result is not
being redone. Spatial/mass generators and finite-G dynamics remain
open. No new external publication is made in this stage.

\begin{thebibliography}{9}
\bibitem{bnr} M. Berkooz, V. Narovlansky and H. Raj,
\emph{Complex Sachdev--Ye--Kitaev model in the double scaling limit},
\href{https://arxiv.org/abs/2006.13983}{arXiv:2006.13983v2},
Section 5, especially (5.2)--(5.5).
\bibitem{tilt} Alan Day and Chengheng, \emph{Joint charge tilt},
local proof, 7 November 2026;
\texttt{susskind-memory-charge-tilt-proof.md}.
\bibitem{joint} Alan Day and Chengheng, \emph{Joint memory regression},
local proof, 7 November 2026, Sections 3--5;
\texttt{susskind-memory-joint-regression-proof.md}.
\bibitem{higher} Alan Day and Chengheng, \emph{One selected state for
higher Wick descendants}, local draft, 8 November 2026.
\bibitem{prep} Alan Day and Chengheng, \emph{Fixed polynomial Slater
preparations}, local draft, 8 November 2026.
\bibitem{belp} L. Ben Efraim and F. Lust-Piquard,
\emph{Poincare type inequalities on the discrete cube and in the CAR algebra},
\href{https://arxiv.org/abs/math/0702233}{arXiv:math/0702233v1},
Section 6.4, Theorem 6.8.
\end{thebibliography}
\end{document}
